% ============================================================================
%  Atlas de Bienestar Humano Territorial — Metodología de construcción
%  Deck para tomadores de decisión · Tensor
%  Compilar:  pdflatex atlas-metodologia.tex  (x2)
% ============================================================================
\documentclass[aspectratio=169,11pt]{beamer}

\usepackage[utf8]{inputenc}
\usepackage[T1]{fontenc}
\usepackage[spanish,es-noshorthands]{babel}
\usepackage{helvet}
\renewcommand{\familydefault}{\sfdefault}
\usepackage{tikz}
\usetikzlibrary{positioning,arrows.meta,shapes.geometric,calc,backgrounds}
\usepackage{booktabs}
\usepackage{array}
\usepackage{fontawesome5}

% ---- Paleta Tensor (modo claro) -------------------------------------------
\definecolor{tPrimary}{HTML}{1A1A1A}
\definecolor{tBody}{HTML}{3D3D3D}
\definecolor{tMuted}{HTML}{5F5F5B}
\definecolor{tAccent}{HTML}{1B6B6D}
\definecolor{tAccentL}{HTML}{E8F4F4}
\definecolor{tAccentD}{HTML}{155556}
\definecolor{tBorder}{HTML}{E5E5E0}
\definecolor{tBgAlt}{HTML}{F2F1EE}
\definecolor{tAmber}{HTML}{B45309}

% ---- Tema minimal ----------------------------------------------------------
\setbeamertemplate{navigation symbols}{}
\setbeamercolor{normal text}{fg=tBody}
\setbeamercolor{structure}{fg=tAccent}
\setbeamercolor{frametitle}{fg=tPrimary}
\setbeamercolor{itemize item}{fg=tAccent}
\setbeamercolor{itemize subitem}{fg=tAccent}
\setbeamercolor{block title}{fg=white,bg=tAccent}
\setbeamercolor{block body}{fg=tBody,bg=tAccentL}
\setbeamerfont{frametitle}{series=\bfseries,size=\Large}
\setbeamerfont{title}{series=\bfseries}
\setbeamertemplate{itemize item}{\footnotesize\faCircle}
\setbeamertemplate{itemize subitem}{\scriptsize\faAngleRight}

% Título de frame con regla accent
\setbeamertemplate{frametitle}{%
  \vspace*{0.7em}%
  {\usebeamerfont{frametitle}\usebeamercolor[fg]{frametitle}\insertframetitle}\\[0.15em]%
  {\color{tAccent}\rule{0.10\textwidth}{2.2pt}}\par\vspace*{0.2em}%
}

% Pie de página
\setbeamertemplate{footline}{%
  \hbox{\begin{beamercolorbox}[wd=\paperwidth,ht=2.6ex,dp=1.2ex,leftskip=1.2em,rightskip=1.2em]{}%
    {\color{tMuted}\scriptsize\texttt{TENSOR}\quad\textbar\quad Atlas · Metodología}%
    \hfill{\color{tMuted}\scriptsize\insertframenumber}%
  \end{beamercolorbox}}%
}

% --- Comandos util ----------------------------------------------------------
\newcommand{\skilltag}[1]{{\color{tAccentD}\ttfamily\footnotesize[#1]}}
\newcommand{\lead}[1]{{\color{tPrimary}\large #1}\par\medskip}
\newcommand{\kpi}[2]{{\color{tAccent}\bfseries\Large #1}\;{\color{tMuted}\footnotesize #2}}

% Caja "Entra / Hace / Sale"
\newcommand{\ehs}[3]{%
\begin{columns}[T]
  \begin{column}{0.32\textwidth}
    {\color{tMuted}\ttfamily\footnotesize ENTRA}\\[2pt]{\small #1}
  \end{column}
  \begin{column}{0.36\textwidth}
    {\color{tMuted}\ttfamily\footnotesize HACE}\\[2pt]{\small #2}
  \end{column}
  \begin{column}{0.32\textwidth}
    {\color{tMuted}\ttfamily\footnotesize PRODUCE}\\[2pt]{\small #3}
  \end{column}
\end{columns}}

% ============================================================================
\begin{document}

% ---------- PORTADA ----------
{\setbeamertemplate{footline}{}
\begin{frame}[plain]
\begin{tikzpicture}[remember picture,overlay]
  \fill[tBgAlt] (current page.south west) rectangle (current page.north east);
  \fill[tAccent] (current page.north west) rectangle ([yshift=-0.9cm]current page.north east);
\end{tikzpicture}
\vspace{1.6cm}
{\color{tAccent}\ttfamily\small ATLAS DE BIENESTAR HUMANO TERRITORIAL}\\[0.5em]
{\color{tPrimary}\fontsize{30}{34}\selectfont\bfseries Cómo se construye un Atlas,\\ paso a paso}\\[0.6em]
{\color{tBody}\large Una metodología repetible para convertir datos públicos\\ en decisiones territoriales — explicada sin tecnicismos.}\\[1.4cm]
{\color{tMuted}\ttfamily\footnotesize TENSOR \quad\textbar\quad uraba.tensor.lat \;\textbullet\; caba.tensor.lat}
\end{frame}}

% ---------- 1. EL PROBLEMA ----------
\begin{frame}{El punto de partida}
\lead{Los datos de un territorio existen — pero están dispersos, en formatos distintos y sin una lectura común.}
\begin{columns}[T]
\begin{column}{0.5\textwidth}
\textbf{\color{tPrimary}Lo que hay hoy}
\begin{itemize}\footnotesize
  \item Catastro en una oficina, censo en otra, salud en otra.
  \item Formatos incompatibles, sin geografía común.
  \item Cada estudio se arma desde cero y no se reutiliza.
\end{itemize}
\end{column}
\begin{column}{0.5\textwidth}
\textbf{\color{tPrimary}La consecuencia}
\begin{itemize}\footnotesize
  \item Decisiones sin granularidad: ``el municipio'', no ``esta manzana''.
  \item Difícil comparar zonas o priorizar con evidencia.
  \item Lo que costó meses no sirve para el siguiente territorio.
\end{itemize}
\end{column}
\end{columns}
\medskip
\begin{block}{La pregunta}
¿Y si midiéramos el \textbf{bienestar de cada cuadra}, lo volviéramos un mapa que cualquiera entienda, y repitiéramos el proceso en cualquier territorio?
\end{block}
\end{frame}

% ---------- 2. QUÉ ES UN ATLAS ----------
\begin{frame}{Qué es un Atlas de Tensor}
\lead{Un mapa interactivo que pinta un \textbf{Índice de Bienestar} sobre miles de unidades pequeñas (manzanas), y deja explorarlo dimensión por dimensión.}
\begin{columns}[T]
\begin{column}{0.55\textwidth}
\textbf{\color{tPrimary}Lo que ve el usuario}
\begin{itemize}
  \item Un mapa donde el color dice qué tan bien está cada zona.
  \item Filtros por dimensión: salud, ambiente, transporte\dots
  \item Fichas por barrio/comuna con sus indicadores.
  \item Comparador de zonas y rankings.
\end{itemize}
\end{column}
\begin{column}{0.45\textwidth}
\begin{block}{Casos reales}
\textbf{Urabá} (Colombia) — 7.028 manzanas.\\[3pt]
\textbf{CABA} (Argentina) — radios censales de Buenos Aires.\\[6pt]
\footnotesize Mismo método, dos países, mismo resultado.
\end{block}
\end{column}
\end{columns}
\end{frame}

% ---------- 3. PARA QUÉ SIRVE ----------
\begin{frame}{Para qué sirve — quién decide mejor con esto}
\begin{columns}[T]
\begin{column}{0.5\textwidth}
\textbf{\color{tAccent}Sector público}
\begin{itemize}\footnotesize
  \item Focalizar inversión social donde más falta.
  \item Priorizar obras de acceso (salud, transporte).
  \item Ordenamiento territorial con evidencia.
\end{itemize}
\medskip
\textbf{\color{tAccent}Gremios y cooperación}
\begin{itemize}\footnotesize
  \item Diagnóstico territorial para proyectos.
  \item Línea base para medir impacto.
\end{itemize}
\end{column}
\begin{column}{0.5\textwidth}
\textbf{\color{tAccent}Sector privado}
\begin{itemize}\footnotesize
  \item Decidir dónde ubicar servicios o expandir.
  \item Entender el entorno socioeconómico de una zona.
\end{itemize}
\medskip
\begin{block}{El valor}
Pasar de la intuición y el promedio municipal a una \textbf{lectura cuadra por cuadra}, comparable y auditable.
\end{block}
\end{column}
\end{columns}
\end{frame}

% ---------- SECCIÓN: CÓMO SE MIDE ----------
\begin{frame}[plain]
\begin{tikzpicture}[remember picture,overlay]
  \fill[tAccent] (current page.south west) rectangle (current page.north east);
\end{tikzpicture}
\vfill
\centering
{\color{white}\ttfamily\small PARTE 1}\\[0.4em]
{\color{white}\Huge\bfseries Cómo se mide el bienestar}\\[0.4em]
{\color{tAccentL}\large El Índice, sus dimensiones y la unidad de análisis}
\vfill
\end{frame}

% ---------- 4. EL ÍNDICE ----------
\begin{frame}{El Índice de Bienestar: cuatro dimensiones}
\lead{El puntaje de cada manzana combina cuatro dimensiones, con pesos fijos. Es transparente: se sabe exactamente de dónde sale cada número.}
\begin{columns}[T]
\begin{column}{0.54\textwidth}
\begin{tikzpicture}[xscale=4.2,yscale=0.62]
  \foreach \y/\name/\w/\val in {3/Accesibilidad/0.40/1.0, 2/Ambiental/0.25/0.625, 1/Socioeconómico/0.25/0.625, 0/Seguridad/0.20/0.50}{
    \node[anchor=east,text=tBody] at (0,\y) {\small \name};
    \fill[tBorder] (0.05,\y-0.28) rectangle (1.05,\y+0.28);
    \fill[tAccent] (0.05,\y-0.28) rectangle (0.05+\val,\y+0.28);
    \node[anchor=west,text=white,font=\scriptsize\bfseries] at (0.10,\y) {\w};
  }
\end{tikzpicture}
\end{column}
\begin{column}{0.46\textwidth}
\footnotesize
\textbf{\color{tPrimary}Cómo se lee}
\begin{itemize}\footnotesize
  \item Cada dimensión vale entre 0 y 1 (peor\,$\rightarrow$\,mejor).
  \item Los pesos reflejan prioridad: el acceso pesa más.
  \item Si falta dato real de una dimensión, el índice se \textbf{re-balancea} con las que sí hay (no se inventa).
\end{itemize}
\smallskip
{\color{tMuted}Índice $=\dfrac{0.40\cdot\text{Acc}+0.25\cdot\text{Amb}+0.25\cdot\text{Soc}+0.20\cdot\text{Seg}}{\text{suma de pesos}}$}
\end{column}
\end{columns}
\end{frame}

% ---------- 5. SUB-INDICADORES ----------
\begin{frame}{Qué hay dentro de cada dimensión}
\lead{Cada dimensión se arma de sub-indicadores medibles. Ejemplos:}
\begin{columns}[T]
\begin{column}{0.5\textwidth}
\footnotesize
\textbf{\color{tAccent}Accesibilidad} — qué tan cerca está lo esencial
\begin{itemize}\footnotesize\item Tiempo real a salud, educación y transporte.\end{itemize}
\smallskip
\textbf{\color{tAccent}Ambiental} — calidad del entorno
\begin{itemize}\footnotesize\item Vegetación (NDVI), islas de calor (temperatura), suelo construido.\end{itemize}
\end{column}
\begin{column}{0.5\textwidth}
\footnotesize
\textbf{\color{tAccent}Socioeconómico} — condiciones de vida
\begin{itemize}\footnotesize\item Necesidades básicas insatisfechas, hogares, población.\end{itemize}
\smallskip
\textbf{\color{tAccent}Seguridad} — proxy disponible
\begin{itemize}\footnotesize\item Iluminación nocturna y presencia de equipamiento.\end{itemize}
\end{column}
\end{columns}
\vfill
\begin{block}{Honestidad metodológica}
Lo que tiene \textbf{dato real} se marca como tal; lo que es aproximación (``proxy'') se declara; lo que falta se registra como \textbf{brecha}. Nada se disfraza.
\end{block}
\end{frame}

% ---------- 6. UNIDAD DE ANÁLISIS ----------
\begin{frame}{La unidad de análisis: del detalle al agregado}
\lead{Medimos en la unidad más fina posible y luego agregamos hacia arriba para los paneles y rankings.}
\centering\medskip
\begin{tikzpicture}[node distance=0.7cm,every node/.style={font=\small}]
  \node[draw=tAccent,fill=tAccentL,rounded corners,thick,minimum width=3.4cm,minimum height=1.1cm,align=center] (m) {\textbf{Manzana / Radio censal}\\[2pt]\scriptsize unidad base del Índice};
  \node[draw=tBorder,fill=white,rounded corners,minimum width=3.0cm,minimum height=1.0cm,align=center,right=1.3cm of m] (b) {\textbf{Barrio / Vereda}\\[2pt]\scriptsize capa de contexto};
  \node[draw=tBorder,fill=white,rounded corners,minimum width=3.0cm,minimum height=1.0cm,align=center,right=1.3cm of b] (mu) {\textbf{Comuna / Municipio}\\[2pt]\scriptsize agregación y rankings};
  \draw[-{Latex[length=3mm]},tAccent,thick] (m) -- (b);
  \draw[-{Latex[length=3mm]},tAccent,thick] (b) -- (mu);
\end{tikzpicture}
\vfill
{\footnotesize\color{tMuted} Miles de manzanas dan el detalle; las agregaciones dan la mirada comparable. Una sola base de datos sirve para las dos escalas.}
\end{frame}

% ---------- 7. DE DÓNDE SALEN LOS DATOS ----------
\begin{frame}{De dónde salen los datos: fuentes abiertas + satélite}
\lead{Todo se construye con datos \textbf{públicos y verificables}, más imágenes satelitales gratuitas.}
\centering\small
\begin{tabular}{@{}>{\raggedright}p{3.3cm} p{8.2cm}@{}}
\toprule
{\color{tAccent}\textbf{Dominio}} & {\color{tAccent}\textbf{Fuentes típicas}}\\
\midrule
Territorio base & Catastro, censo de población (manzanas, NBI, hogares)\\
Salud / Educación & Registros de hospitales, centros de salud, colegios\\
Transporte & Paradas, rutas, ciclovías; tiempos de viaje reales\\
Ambiente & Áreas verdes, arbolado; satélite (vegetación, calor)\\
Riesgo / Social & Asentamientos, inundación, vulnerabilidad\\
\bottomrule
\end{tabular}
\vfill
{\footnotesize\color{tMuted} En Colombia: IGAC, DANE, REPS, SIMAT, IDEAM\dots\quad En Argentina: BA Data, INDEC, IGN\dots\quad Satélite (mundial): Sentinel-2, Landsat, VIIRS.}
\end{frame}

% ---------- SECCIÓN: EL MÉTODO ----------
\begin{frame}[plain]
\begin{tikzpicture}[remember picture,overlay]
  \fill[tAccent] (current page.south west) rectangle (current page.north east);
\end{tikzpicture}
\vfill
\centering
{\color{white}\ttfamily\small PARTE 2}\\[0.4em]
{\color{white}\Huge\bfseries El método de construcción}\\[0.4em]
{\color{tAccentL}\large Seis fases, de los datos crudos al atlas en producción}
\vfill
\end{frame}

% ---------- 8. POR QUÉ ES REPETIBLE ----------
\begin{frame}{La clave: una fábrica, no una obra irrepetible}
\lead{El visor del mapa no sabe de qué territorio habla. Solo exige que los datos lleguen en un \textbf{molde fijo} (un ``contrato'').}
\begin{columns}[T]
\begin{column}{0.5\textwidth}
\begin{block}{Consecuencia}
Construir un atlas nuevo $=$ \textbf{reemplazar los datos} en el molde $+$ cambiar nombres y marca.\\[4pt] No se reescribe el programa.
\end{block}
\end{column}
\begin{column}{0.5\textwidth}
\textbf{\color{tPrimary}Por eso\dots}
\begin{itemize}\footnotesize
  \item El primer atlas costó meses de aprendizaje.
  \item El segundo (CABA) se armó en una fracción del tiempo.
  \item El método queda \textbf{empaquetado} y se aplica a cualquier territorio.
\end{itemize}
\end{column}
\end{columns}
\end{frame}

% ---------- 9. PIPELINE GENERAL ----------
\begin{frame}{Las seis fases de un vistazo}
\centering\medskip
\begin{tikzpicture}[
  fase/.style={draw=tAccent,fill=white,rounded corners,thick,minimum width=2.05cm,minimum height=1.35cm,align=center,font=\scriptsize},
  node distance=0.45cm]
  \node[fase] (f0) {{\color{tAccent}\ttfamily\bfseries 0}\\[1pt]\textbf{Explorar}\\datos};
  \node[fase,right=of f0] (f1) {{\color{tAccent}\ttfamily\bfseries 1}\\[1pt]\textbf{Diseñar}\\el índice};
  \node[fase,right=of f1] (f2) {{\color{tAccent}\ttfamily\bfseries 2}\\[1pt]\textbf{Procesar}\\los datos};
  \node[fase,right=of f2] (f3) {{\color{tAccent}\ttfamily\bfseries 3}\\[1pt]\textbf{Satélite}\\(opcional)};
  \node[fase,right=of f3] (f4) {{\color{tAccent}\ttfamily\bfseries 4}\\[1pt]\textbf{Construir}\\el visor};
  \node[fase,right=of f4] (f5) {{\color{tAccent}\ttfamily\bfseries 5}\\[1pt]\textbf{Publicar}\\y verificar};
  \foreach \a/\b in {f0/f1,f1/f2,f2/f3,f3/f4,f4/f5}{
    \draw[-{Latex[length=2.4mm]},tAccent,thick] (\a) -- (\b);}
\end{tikzpicture}
\vfill
\begin{columns}[T]
\begin{column}{0.62\textwidth}
{\footnotesize\color{tBody} Entre cada fase hay un \textbf{control de calidad} (``compuerta''): no se avanza hasta que la fase anterior cumple su criterio. Si algo falla, se detiene y se reporta.}
\end{column}
\begin{column}{0.34\textwidth}
{\footnotesize\color{tMuted} Cada fase está documentada como una guía ejecutable reutilizable. El proceso completo puede reanudarse donde se quedó.}
\end{column}
\end{columns}
\end{frame}

% ---------- 10-15. FASES ----------
\begin{frame}{Fase 0 — Explorar y verificar las fuentes \skilltag{atlas-scout}}
\lead{Buscar qué datos abiertos existen para el territorio y \textbf{comprobar uno por uno} que estén vivos y traigan información real.}
\ehs{El nombre del territorio y su país.}{Rastrea fuentes por dominio (salud, transporte\dots) y verifica cada una en el momento: ¿responde?, ¿trae datos?, ¿qué licencia?}{Un inventario de fuentes confiables y una lista honesta de \textbf{vacíos} de información.}
\vfill
\begin{block}{Por qué importa}
Evita construir sobre datos que ``parecían'' existir. Lo que no se puede verificar se declara como brecha, no se rellena con supuestos.
\end{block}
\end{frame}

\begin{frame}{Fase 1 — Diseñar el índice \skilltag{atlas-design}}
\lead{Decidir la unidad de análisis, las dimensiones y los pesos según lo que el dato \textbf{permite} en ese territorio.}
\ehs{El inventario de fuentes de la Fase 0.}{Elige la unidad base y las agregaciones; fija dimensiones y pesos; clasifica cada una como dato-real, proxy o brecha.}{Un plan de construcción claro, acordado antes de procesar nada.}
\vfill
\begin{block}{Compuerta de calidad}
El plan se \textbf{confirma con el responsable} antes de seguir. Aquí se acuerda qué se va a medir y cómo.
\end{block}
\end{frame}

\begin{frame}{Fase 2 — Procesar los datos \skilltag{atlas-pipeline}}
\lead{Convertir los datos crudos en el ``molde'' que el visor entiende, calculando el Índice de cada manzana.}
\ehs{Datos crudos verificados + el plan.}{Limpia, alinea geografías, normaliza, calcula los puntajes ponderados y prepara los mapas para que carguen rápido.}{La base de datos del atlas, lista y \textbf{auto-verificada}.}
\vfill
\begin{block}{Compuerta de calidad}
Verificación automática: todos los puntajes entre 0 y 1, sin valores vacíos, y los conteos cuadran. Si no pasa, no avanza.
\end{block}
\end{frame}

\begin{frame}{Fase 3 — Indicadores satelitales \skilltag{atlas-satelite}\;\;\small(opcional)}
\lead{Medir el entorno físico con imágenes de satélite gratuitas, manzana por manzana.}
\ehs{La geometría de las manzanas.}{Calcula vegetación, temperatura de superficie (islas de calor), suelo construido y luces nocturnas sobre cada unidad.}{Una tabla de indicadores ambientales que enriquece el Índice.}
\vfill
{\footnotesize\color{tMuted} Usa programas satelitales abiertos y globales (Sentinel-2, Landsat, VIIRS), así que aplica a cualquier territorio del mundo sin pedir permisos.}
\end{frame}

\begin{frame}{Fase 4 — Construir el visor \skilltag{atlas-web}}
\lead{Tomar el visor de mapa ya probado y \textbf{vestirlo} con los datos y la identidad del nuevo territorio.}
\ehs{La base de datos de la Fase 2.}{Reusa el mismo visor; cambia nombres, etiquetas, centro del mapa y marca. No se reescribe el programa.}{El atlas funcionando, pintando el nuevo territorio.}
\vfill
\begin{block}{Eficiencia}
Como el visor es el mismo molde, esta fase es sobre todo configuración — el grueso del trabajo ya quedó en los datos.
\end{block}
\end{frame}

\begin{frame}{Fase 5 — Publicar y verificar \skilltag{atlas-deploy}}
\lead{Poner el atlas en línea, con su dirección web, tras un control de calidad final.}
\ehs{El atlas funcionando localmente.}{Publica en la nube, asigna el dominio y revisa \textbf{dato por dato y pantalla por pantalla} que todo sea coherente.}{Un atlas público, en una dirección web, listo para decidir.}
\vfill
{\footnotesize\color{tMuted} Doble control: que los números sean consistentes y que el mapa y los paneles se vean y funcionen bien.}
\end{frame}

% ---------- 16. CONTROL DE CALIDAD ----------
\begin{frame}{El hilo conductor: control de calidad entre fases}
\lead{Lo que hace confiable al método no es una fase, sino las \textbf{compuertas} entre ellas.}
\begin{columns}[T]
\begin{column}{0.5\textwidth}
\textbf{\color{tPrimary}Principios}
\begin{itemize}\footnotesize
  \item No se avanza sin verificar la fase anterior.
  \item Toda fuente se comprueba en el momento.
  \item Lo no verificable se declara como brecha.
  \item El proceso es reanudable y auditable.
\end{itemize}
\end{column}
\begin{column}{0.5\textwidth}
\begin{block}{Resultado}
Un atlas que se puede \textbf{defender}: cada número tiene origen, cada aproximación está declarada, cada vacío está documentado.
\end{block}
\end{column}
\end{columns}
\end{frame}

% ---------- 17. AUTOMATIZACIÓN ----------
\begin{frame}{Cómo se ejecuta este método}
\lead{Cada fase está escrita como una \textbf{guía ejecutable} que un sistema de agentes de inteligencia artificial sigue de forma asistida y supervisada.}
\begin{columns}[T]
\begin{column}{0.5\textwidth}
\textbf{\color{tPrimary}Qué significa en la práctica}
\begin{itemize}\footnotesize
  \item El método queda \textbf{empaquetado}, no en la cabeza de una persona.
  \item Tareas repetitivas (buscar, verificar, procesar) se automatizan.
  \item Una persona supervisa y aprueba en las compuertas.
\end{itemize}
\end{column}
\begin{column}{0.5\textwidth}
\begin{block}{Lo importante para el decisor}
No necesita conocer la herramienta. Lo relevante es que el método es \textbf{repetible, rápido y auditable} — y que existe el mismo para cualquier territorio.
\end{block}
\end{column}
\end{columns}
\end{frame}

% ---------- 18. EL SALTO ----------
\begin{frame}{El salto: de meses a días}
\centering\vspace{0.3cm}
\begin{columns}[T]
\begin{column}{0.33\textwidth}\centering
\kpi{Meses}{primer atlas}\\[4pt]{\footnotesize\color{tMuted} Urabá: definir el método desde cero}
\end{column}
\begin{column}{0.33\textwidth}\centering
\kpi{Días}{atlas siguiente}\\[4pt]{\footnotesize\color{tMuted} CABA: aplicar el método ya probado}
\end{column}
\begin{column}{0.33\textwidth}\centering
\kpi{$N$}{territorios}\\[4pt]{\footnotesize\color{tMuted} El mismo molde escala a cualquiera}
\end{column}
\end{columns}
\vfill
\begin{block}{En palabras del propio análisis}
Lo que en un territorio costó meses de trabajo, en el siguiente quedó ``a un clic de distancia'' — porque las piezas difíciles ya están resueltas y empaquetadas.
\end{block}
\end{frame}

% ---------- 19. VALOR / CIERRE ----------
\begin{frame}{Qué gana un tomador de decisiones}
\begin{columns}[T]
\begin{column}{0.5\textwidth}
\textbf{\color{tAccent}Decisión con evidencia}
\begin{itemize}\footnotesize
  \item Ver el bienestar cuadra por cuadra.
  \item Priorizar dónde actuar primero.
  \item Comparar zonas con un criterio único.
\end{itemize}
\medskip
\textbf{\color{tAccent}Confianza}
\begin{itemize}\footnotesize
  \item Datos públicos y verificables.
  \item Método transparente y auditable.
\end{itemize}
\end{column}
\begin{column}{0.5\textwidth}
\textbf{\color{tAccent}Escala}
\begin{itemize}\footnotesize
  \item Un territorio hoy, varios mañana.
  \item Línea base para medir cambios en el tiempo.
\end{itemize}
\medskip
\begin{block}{En síntesis}
Un instrumento para gobernar y invertir \textbf{con el territorio a la vista} — no con promedios.
\end{block}
\end{column}
\end{columns}
\end{frame}

% ---------- CIERRE ----------
{\setbeamertemplate{footline}{}
\begin{frame}[plain]
\begin{tikzpicture}[remember picture,overlay]
  \fill[tBgAlt] (current page.south west) rectangle (current page.north east);
  \fill[tAccent] ([yshift=0.9cm]current page.south west) rectangle (current page.south east);
\end{tikzpicture}
\vspace{1.8cm}
{\color{tPrimary}\Huge\bfseries Del dato disperso\\ a la decisión territorial.}\\[0.8em]
{\color{tBody}\large La metodología, empaquetada y repetible, para construir un Atlas en cualquier territorio.}\\[1.6cm]
{\color{tAccentD}\ttfamily\small TENSOR}\quad{\color{tMuted}\ttfamily\small\textbar\quad uraba.tensor.lat \textbullet\ caba.tensor.lat \textbullet\ github.com/Cespial/tensor-atlas}
\end{frame}}

\end{document}
